\section{Konu 6: Zincir Kuralı, Yönlü Türevler ve Gradyan Vektörü}

Bu bölümde derste (5 ve 8 Kasım 2026) tahtada çözülen tüm sorular, işlenme sırasıyla ve adım adım çözümleriyle yer alıyor.
Yönlü türev, gradyan ve teğet düzlem konuları derste tahtada işlenmediği için bunlara ait sorular sondaki ek
alıştırmalarda verilmiştir.

\begin{bilgi}[Kullanılan temel formüller]
\begin{itemize}
  \item $z=f(x,y)$, $x(t),y(t)$: $\dfrac{dz}{dt}=z_x\dfrac{dx}{dt}+z_y\dfrac{dy}{dt}$.
  \item $z=f(x,y)$, $x(u,v),y(u,v)$: $z_u=z_xx_u+z_yy_u$, $z_v=z_xx_v+z_yy_v$ (ağaçta $u$'ya inen tüm yollar).
  \item Kapalı türev: $F(x,y)=0\Rightarrow y'=-F_x/F_y$; $F(x,y,z)=0\Rightarrow z_x=-F_x/F_z$, $z_y=-F_y/F_z$.
  \item $D_{\mathbf u}f=\nabla f\cdot\mathbf u$ ($|\mathbf u|=1$); en büyük yönlü türev $|\nabla f|$, $\nabla f$ yönünde.
\end{itemize}
\end{bilgi}

\subsection{5 Kasım dersi: zincir kuralı problemleri}

\begin{soru}[Soru 1 (tahtada Problem 6) \foto{IMG-20261105-WA0041}, \foto{IMG-20261108-WA0087}]
$\dfrac{dz}{dt}$'yi (i) zincir kuralıyla, (ii) $z$'yi $t$ cinsinden yazıp doğrudan türev alarak bulunuz.\\
\textbf{a)} $z=x^2+y^2,\quad x=\cos t+\sin t,\quad y=\cos t-\sin t$\\
\textbf{b)} $z=\arctan\dfrac{y}{x},\quad x=\cos t,\quad y=\sin t$
\end{soru}
\begin{cevap}[Çözüm (adım adım)]
\textbf{a) (i) Zincir kuralı.} $z$, $x$ ve $y$ üzerinden yalnızca $t$'ye bağlı olduğundan Durum 1 uygulanır:
$\frac{dz}{dt}=z_x\frac{dx}{dt}+z_y\frac{dy}{dt}$.
\begin{itemize}
  \item Dış türevler: $z_x=2x$, $z_y=2y$.
  \item İç türevler: $x'=-\sin t+\cos t$, $y'=-\sin t-\cos t$.
\end{itemize}
\begin{align*}
\frac{dz}{dt}&=2(\cos t+\sin t)(\cos t-\sin t)+2(\cos t-\sin t)(-\sin t-\cos t)\\
&=2(\cos^2t-\sin^2t)-2(\cos t-\sin t)(\cos t+\sin t)\\
&=2(\cos^2t-\sin^2t)-2(\cos^2t-\sin^2t)=\boxed{0}.
\end{align*}
\textbf{(ii) Doğrudan.} $(\cos t\pm\sin t)^2=1\pm2\sin t\cos t$ olduğundan $z=(1+2\sin t\cos t)+(1-2\sin t\cos t)=2$.
Sabit fonksiyonun türevi $0$'dır: $dz/dt=0$. İki yol uyuşuyor.

\textbf{b) (i) Zincir kuralı.} $\dfrac{d}{dw}\arctan w=\dfrac{1}{1+w^2}$ ve $w=y/x$ ile:
\[
z_x=\frac{1}{1+\frac{y^2}{x^2}}\cdot\left(-\frac{y}{x^2}\right)=-\frac{y}{x^2+y^2},\qquad
z_y=\frac{1}{1+\frac{y^2}{x^2}}\cdot\frac{1}{x}=\frac{x}{x^2+y^2}.
\]
$x=\cos t$, $y=\sin t$ için $x^2+y^2=1$, $x'=-\sin t$, $y'=\cos t$:
\[
\frac{dz}{dt}=\underbrace{(-\sin t)}_{z_x}\underbrace{(-\sin t)}_{x'}+\underbrace{\cos t}_{z_y}\underbrace{\cos t}_{y'}
=\sin^2t+\cos^2t=\boxed{1}.
\]
\textbf{(ii) Doğrudan.} $z=\arctan\dfrac{\sin t}{\cos t}=\arctan(\tan t)=t$ ($-\frac\pi2<t<\frac\pi2$), dolayısıyla
$dz/dt=1$. (Tahtada b şıkkının (ii) yolu yazılmamıştı; kaynaklara göre tamamlandı.)
\end{cevap}

\begin{soru}[Soru 2 (tahtada Problem 7) \foto{IMG-20261105-WA0042}]
$x=e^t$ dönüşümü ile $x^2\dfrac{d^2y}{dx^2}+x\dfrac{dy}{dx}+y=0$ denklemini $t$ değişkenine göre yazınız.
\end{soru}
\begin{cevap}[Çözüm (adım adım)]
\textbf{1. Dönüşümün türevi.} $x=e^t\iff t=\ln x$ ($x>0$); $\dfrac{dt}{dx}=\dfrac1x=e^{-t}$.

\textbf{2. Birinci türev (zincir kuralı).} $y$, $t$ üzerinden $x$'e bağlıdır:
$\dfrac{dy}{dx}=\dfrac{dy}{dt}\cdot\dfrac{dt}{dx}=e^{-t}\dfrac{dy}{dt}$.

\textbf{3. İkinci türev.} $\frac{dy}{dx}=e^{-t}\frac{dy}{dt}$ bir $t$ fonksiyonudur; $x$'e göre türevi için önce $t$'ye göre
türevleyip (çarpım kuralı) $\frac{dt}{dx}$ ile çarparız:
\begin{align*}
\frac{d^2y}{dx^2}&=\frac{d}{dt}\left(e^{-t}\frac{dy}{dt}\right)\cdot\frac{dt}{dx}
=\left(-e^{-t}\frac{dy}{dt}+e^{-t}\frac{d^2y}{dt^2}\right)e^{-t}
=e^{-2t}\left(\frac{d^2y}{dt^2}-\frac{dy}{dt}\right).
\end{align*}
\textbf{4. Katsayılarla çarpım.} $x^2=e^{2t}$, $x=e^t$:
\[
x^2\frac{d^2y}{dx^2}=\frac{d^2y}{dt^2}-\frac{dy}{dt},\qquad x\frac{dy}{dx}=\frac{dy}{dt}.
\]
\textbf{5. Denklem.}
$\dfrac{d^2y}{dt^2}-\dfrac{dy}{dt}+\dfrac{dy}{dt}+y=0\ \Longrightarrow\ \boxed{\dfrac{d^2y}{dt^2}+y=0}$.

Bu sabit katsayılı denklemin genel çözümü $y=c_1\cos t+c_2\sin t$, yani $y=c_1\cos(\ln x)+c_2\sin(\ln x)$'tir.

\textbf{Not:} Tahtada son satırda ``$x\cdot\frac{dx}{dy}$'' yazılmıştı; doğrusu $x\cdot\frac{dy}{dx}$'tir. Sonuç denklemi
fotoğrafta kadraj dışında kaldı \okunamadi{sağ kenardaki kutulu sonuç}; kaynaklara göre doğrusu $\frac{d^2y}{dt^2}+y=0$'dır.
\end{cevap}

\begin{soru}[Soru 3 (tahtada Problem 8, kısmen görünüyor) \foto{IMG-20261105-WA0042}]
\okunamadi{sorunun metni ve çözümü fotoğrafın sağ kenarında kesik; yalnızca ``8) a, \dots'', ``$x-a\dots$'',
``oldu\dots'', ``$F(x\dots$'', ``$F_x$'', ``$\partial z/\partial x$'', ``$\frac{a}{a-b\dots}$'' parçaları seçilebiliyor}.
Görünen parçalar kaynaklarda yaygın olan şu soruyla uyumludur (\textbf{tahmini metin}):
$a,b$ sabit, $\varphi$ türevlenebilir bir fonksiyon ve $x-az=\varphi(y-bz)$ ise $a\dfrac{\partial z}{\partial x}+b\dfrac{\partial z}{\partial y}=1$
olduğunu gösteriniz.
\end{soru}
\begin{cevap}[Çözüm (adım adım)]
\textbf{1. Kapalı biçim.} $F(x,y,z)=x-az-\varphi(y-bz)=0$. Burada $\varphi'$, $w=y-bz$ noktasında hesaplanır.

\textbf{2. Kısmi türevler} (Durum 4: $\varphi$ tek değişkenli, $w$ iki değişkenli):
\[
F_x=1,\qquad F_y=-\varphi'(w)\cdot w_y=-\varphi',\qquad F_z=-a-\varphi'(w)\cdot w_z=-a+b\varphi'=-(a-b\varphi').
\]
\textbf{3. Kapalı türev formülleri} ($a-b\varphi'\ne0$):
\[
\frac{\partial z}{\partial x}=-\frac{F_x}{F_z}=\frac{1}{a-b\varphi'},\qquad
\frac{\partial z}{\partial y}=-\frac{F_y}{F_z}=-\frac{\varphi'}{a-b\varphi'}.
\]
\textbf{4. İstenen ifade.}
\[
a\frac{\partial z}{\partial x}+b\frac{\partial z}{\partial y}=\frac{a}{a-b\varphi'}-\frac{b\varphi'}{a-b\varphi'}=\frac{a-b\varphi'}{a-b\varphi'}=1 .
\]
Tahtada görünen ``$\frac{a}{a-b\dots}$'' parçası, $a\,\partial z/\partial x$ teriminin bu biçimine karşılık geliyor.
\end{cevap}

\subsection{8 Kasım dersi: zincir kuralı durumları ve kapalı türev}

\begin{soru}[Soru 4 \foto{IMG-20261108-WA0072}, \foto{IMG-20261108-WA0073}, \foto{IMG-20261108-WA0074}]
$z=\ln(x^2+y^2)$, $x=e^{v}\cos u$, $y=e^{v}\sin u$ ise $\dfrac{\partial z}{\partial u}$ ve $\dfrac{\partial z}{\partial v}$'yi bulunuz.
\end{soru}
\begin{cevap}[Çözüm (adım adım)]
\textbf{1. Ağaç.} $z\to x,y$; $x\to u,v$; $y\to u,v$. Formüller:
$z_u=z_xx_u+z_yy_u$, $\ z_v=z_xx_v+z_yy_v$.

\textbf{2. Dış türevler.} $z_x=\dfrac{2x}{x^2+y^2}$, $z_y=\dfrac{2y}{x^2+y^2}$; ayrıca $x^2+y^2=e^{2v}(\cos^2u+\sin^2u)=e^{2v}$.

\textbf{3. İç türevler.} $x_u=-e^{v}\sin u$, $x_v=e^{v}\cos u$, $y_u=e^{v}\cos u$, $y_v=e^{v}\sin u$.

\textbf{4. Hesap.}
\begin{align*}
\frac{\partial z}{\partial u}&=\frac{2e^{v}\cos u\,(-e^{v}\sin u)+2e^{v}\sin u\;e^{v}\cos u}{e^{2v}}=\frac{-2e^{2v}\sin u\cos u+2e^{2v}\sin u\cos u}{e^{2v}}=\boxed{0},\\[4pt]
\frac{\partial z}{\partial v}&=\frac{2e^{v}\cos u\;e^{v}\cos u+2e^{v}\sin u\;e^{v}\sin u}{e^{2v}}=\frac{2e^{2v}(\cos^2u+\sin^2u)}{e^{2v}}=\boxed{2}.
\end{align*}
\textbf{5. Kontrol.} $z=\ln(e^{2v})=2v$ $\Rightarrow$ $z_u=0$, $z_v=2$. \checkmark
(Tahtadaki kutuda üs ``$e^{2u}$'' gibi okunuyor; doğrusu $e^{2v}$.)
\end{cevap}

\begin{soru}[Soru 5 \foto{IMG-20261108-WA0075}]
$z=\dfrac{1}{x^2+y^2}$, $x=\cos2t$, $y=\sin2t$ ise $\dfrac{dz}{dt}$'yi bulunuz.
\end{soru}
\begin{cevap}[Çözüm (adım adım)]
\textbf{1.} $x^2+y^2=\cos^22t+\sin^22t=1$.

\textbf{2. Kısmi türevler.} $z=(x^2+y^2)^{-1}$ olduğundan
$z_x=-(x^2+y^2)^{-2}\cdot2x=\dfrac{-2x}{(x^2+y^2)^2}$, $z_y=\dfrac{-2y}{(x^2+y^2)^2}$.

\textbf{3. İç türevler.} $x'=-2\sin2t$, $y'=2\cos2t$ (içteki $2t$'nin türevi $2$ unutulmamalı).

\textbf{4. Durum 1.}
\[
\frac{dz}{dt}=(-2\cos2t)(-2\sin2t)+(-2\sin2t)(2\cos2t)=4\sin2t\cos2t-4\sin2t\cos2t=\boxed{0}.
\]
\textbf{5. Kontrol.} $z=\dfrac{1}{\cos^22t+\sin^22t}=1\Rightarrow dz/dt=0$. (Tahtada ``$z=$'' yazılıp bırakılmıştı.)
\end{cevap}

\begin{soru}[Soru 6 \foto{IMG-20261108-WA0077}, \foto{IMG-20261108-WA0078}]
$s=\ln(x^2+y^2+z^2)$, $x=u-v$, $y=u+v$, $z=2\sqrt{uv}$ ise $\dfrac{\partial s}{\partial u}$ ve $\dfrac{\partial s}{\partial v}$'yi bulunuz.
\end{soru}
\begin{cevap}[Çözüm (adım adım)]
\textbf{1. Ağaç (Durum 3).} $s\to x,y,z$; her biri $\to u,v$. Her türevde üç terim:
$s_u=s_xx_u+s_yy_u+s_zz_u$, $\ s_v=s_xx_v+s_yy_v+s_zz_v$.

\textbf{2. Dış türevler.} $s_x=\dfrac{2x}{N}$, $s_y=\dfrac{2y}{N}$, $s_z=\dfrac{2z}{N}$, burada
$N=x^2+y^2+z^2=(u-v)^2+(u+v)^2+4uv=2u^2+2v^2+4uv=2(u+v)^2$.

\textbf{3. İç türevler.} $x_u=1,\ x_v=-1,\ y_u=1,\ y_v=1,\ z_u=\dfrac{v}{\sqrt{uv}},\ z_v=\dfrac{u}{\sqrt{uv}}$
($z=2\sqrt{uv}$ için $z_u=2\cdot\frac{v}{2\sqrt{uv}}$).

\textbf{4. Hesap.} $2z\,z_u=2\cdot2\sqrt{uv}\cdot\frac{v}{\sqrt{uv}}=4v$ ve $2z\,z_v=4u$ olduğundan
\begin{align*}
\frac{\partial s}{\partial u}&=\frac{2(u-v)+2(u+v)+4v}{2(u+v)^2}=\frac{4u+4v}{2(u+v)^2}=\boxed{\frac{2}{u+v}},\\
\frac{\partial s}{\partial v}&=\frac{-2(u-v)+2(u+v)+4u}{2(u+v)^2}=\frac{4u+4v}{2(u+v)^2}=\boxed{\frac{2}{u+v}}.
\end{align*}
\textbf{5. Kontrol.} $s=\ln\big(2(u+v)^2\big)=\ln2+2\ln(u+v)\Rightarrow s_u=s_v=\dfrac{2}{u+v}$. \checkmark
\end{cevap}

\begin{soru}[Soru 7 \foto{IMG-20261108-WA0079}]
$y=\sin x$, $x=u^2v$ ise $y_u$ ve $y_v$'yi bulunuz.
\end{soru}
\begin{cevap}[Çözüm (adım adım)]
\textbf{1. Ağaç (Durum 4).} $y\to x$ (adi türev $dy/dx$), $x\to u,v$ (kısmi türevler).

\textbf{2.} $\dfrac{dy}{dx}=\cos x=\cos(u^2v)$; $x_u=2uv$, $x_v=u^2$.

\textbf{3.} $y_u=\dfrac{dy}{dx}x_u=\boxed{2uv\cos(u^2v)}$,\qquad $y_v=\dfrac{dy}{dx}x_v=\boxed{u^2\cos(u^2v)}$.
\end{cevap}

\begin{soru}[Soru 8 \foto{IMG-20261108-WA0080}, \foto{IMG-20261108-WA0081}]
$2x^4z^3+3xy^2+4xz+3y+5x-23=0$ denklemi $z=z(x,y)$ fonksiyonunu tanımlıyor. $(1,2)$ noktasında $z_x$ ve $z_y$'yi bulunuz.
\end{soru}
\begin{cevap}[Çözüm (adım adım)]
\textbf{1. Formül (WA0080).} $F(x,y,z)=0$'ın $x$'e göre türevi: $F_x\cdot1+F_y\cdot0+F_zz_x=0\Rightarrow z_x=-F_x/F_z$;
benzer şekilde $z_y=-F_y/F_z$ ($F_z\ne0$).

\textbf{2. Noktayı tamamla.} $x=1,\ y=2$: $2z^3+12+4z+6+5-23=0\Rightarrow2z^3+4z=0\Rightarrow2z(z^2+2)=0$.
$z^2+2>0$ olduğundan tek reel çözüm $z=0$; nokta $(1,2,0)$.

\textbf{3. Kısmi türevler} (değişkenler bağımsızmış gibi):
$F_x=8x^3z^3+3y^2+4z+5$, $F_y=6xy+3$, $F_z=6x^4z^2+4x$.

\textbf{4. Noktada:} $F_x=17$, $F_y=15$, $F_z=4\ne0$.

\textbf{5. Sonuç:} $z_x=\boxed{-\dfrac{17}{4}}$, $\quad z_y=\boxed{-\dfrac{15}{4}}$.
\end{cevap}

\subsection{Ek alıştırmalar}

\begin{soru}[\ekal Ek alıştırma 1 (Durum 2)]
$w=xy+z$, $x=\cos t$, $y=\sin t$, $z=t$ ise $dw/dt$'yi bulunuz ve $t=0$'daki değerini hesaplayınız.
\end{soru}
\begin{cevap}
$w_x=y,\ w_y=x,\ w_z=1$; $x'=-\sin t,\ y'=\cos t,\ z'=1$.
\[
\frac{dw}{dt}=\sin t(-\sin t)+\cos t\cos t+1=\cos^2t-\sin^2t+1=\cos2t+1 .
\]
$t=0$'da $dw/dt=2$. Kontrol: $w=\sin t\cos t+t=\frac12\sin2t+t$, $w'=\cos2t+1$. \checkmark
\end{cevap}

\begin{soru}[\ekal Ek alıştırma 2 (kutupsal koordinatlar)]
$z=x^2+xy$, $x=r\cos\theta$, $y=r\sin\theta$ ise $z_r$ ve $z_\theta$'yı zincir kuralıyla bulunuz.
\end{soru}
\begin{cevap}
$z_x=2x+y$, $z_y=x$; $x_r=\cos\theta$, $y_r=\sin\theta$, $x_\theta=-r\sin\theta$, $y_\theta=r\cos\theta$.
\begin{align*}
z_r&=(2x+y)\cos\theta+x\sin\theta=(2r\cos\theta+r\sin\theta)\cos\theta+r\cos\theta\sin\theta
=2r\cos^2\theta+r\sin2\theta,\\
z_\theta&=-(2x+y)\,r\sin\theta+x\,r\cos\theta=-2r^2\sin\theta\cos\theta-r^2\sin^2\theta+r^2\cos^2\theta
=r^2(\cos2\theta-\sin2\theta).
\end{align*}
Kontrol: $z=r^2\cos^2\theta+r^2\sin\theta\cos\theta=r^2\left(\cos^2\theta+\tfrac12\sin2\theta\right)$; doğrudan türevler
aynı sonucu verir. \checkmark
\end{cevap}

\begin{soru}[\ekal Ek alıştırma 3 (Durum 4, keyfi fonksiyon)]
$f$ türevlenebilir ve $z=f(x^2-y^2)$ ise $y\,z_x+x\,z_y=0$ olduğunu gösteriniz.
\end{soru}
\begin{cevap}
$w=x^2-y^2$ dersek $z=f(w)$ (Durum 4): $z_x=f'(w)\,w_x=2x\,f'(w)$, $z_y=f'(w)\,w_y=-2y\,f'(w)$. Buna göre
$y\,z_x+x\,z_y=2xy\,f'-2xy\,f'=0$.
\end{cevap}

\begin{soru}[\ekal Ek alıştırma 4 (kapalı türev)]
$x^3+y^3+z^3+6xyz=1$ denklemi $z=z(x,y)$'yi tanımlıyorsa $z_x$ ve $z_y$'yi bulunuz.
\end{soru}
\begin{cevap}
$F=x^3+y^3+z^3+6xyz-1$; $F_x=3x^2+6yz$, $F_y=3y^2+6xz$, $F_z=3z^2+6xy$. $F_z\ne0$ iken
\[
z_x=-\frac{3x^2+6yz}{3z^2+6xy}=-\frac{x^2+2yz}{z^2+2xy},\qquad z_y=-\frac{y^2+2xz}{z^2+2xy}.
\]
\end{cevap}

\begin{soru}[\ekal Ek alıştırma 5 (Euler--Cauchy)]
$x=e^t$ dönüşümüyle $x^2y''+3xy'+y=0$ denklemini $t$'ye göre yazınız ve çözünüz.
\end{soru}
\begin{cevap}
$x^2y''=\ddot y-\dot y$ ve $xy'=\dot y$ (nokta $t$'ye göre türev) kullanılırsa
\[
\ddot y-\dot y+3\dot y+y=0\ \Longrightarrow\ \ddot y+2\dot y+y=0 .
\]
Karakteristik denklem $m^2+2m+1=(m+1)^2=0$, çift kök $m=-1$: $y=(c_1+c_2t)e^{-t}$. $t=\ln x$ ile
$y=\dfrac{c_1+c_2\ln x}{x}$. (Kontrol: $y=1/x$ için $x^2\cdot\frac{2}{x^3}+3x\cdot\left(-\frac{1}{x^2}\right)+\frac1x=0$.)
\end{cevap}

\begin{soru}[\ekal Ek alıştırma 6 (yönlü türev)]
$f(x,y)=x^2y^3-4y$ fonksiyonunun $(2,-1)$ noktasında $\mathbf v=2\mathbf i+5\mathbf j$ yönündeki yönlü türevini bulunuz.
\end{soru}
\begin{cevap}
$f_x=2xy^3$, $f_y=3x^2y^2-4$; $(2,-1)$'de $f_x=2\cdot2\cdot(-1)=-4$, $f_y=3\cdot4\cdot1-4=8$; $\nabla f=\langle-4,8\rangle$.
$\mathbf v$ birim değil: $|\mathbf v|=\sqrt{29}$, $\mathbf u=\frac{1}{\sqrt{29}}\langle2,5\rangle$.
\[
D_{\mathbf u}f(2,-1)=\frac{-4\cdot2+8\cdot5}{\sqrt{29}}=\boxed{\frac{32}{\sqrt{29}}}\approx5{,}94 .
\]
\end{cevap}

\begin{soru}[\ekal Ek alıştırma 7 (en hızlı artış, verilen noktaya doğru yön)]
$f(x,y)=x\,e^{y}$ ve $P(2,0)$ olsun. (a) $P$'den $Q(\frac12,2)$'ye doğru yöndeki yönlü türevi bulunuz. (b) $f$ en hızlı
hangi yönde artar, en büyük artış hızı nedir? (c) Hangi yönlerde değişim hızı sıfırdır?
\end{soru}
\begin{cevap}
$\nabla f=\langle e^{y},\ x e^{y}\rangle$, $\nabla f(2,0)=\langle1,2\rangle$.

\textbf{(a)} $\overrightarrow{PQ}=\langle-\frac32,2\rangle$, $|\overrightarrow{PQ}|=\frac52$, $\mathbf u=\langle-\frac35,\frac45\rangle$;
$D_{\mathbf u}f=-\frac35+\frac85=\boxed{1}$.

\textbf{(b)} $\nabla f(2,0)=\langle1,2\rangle$ yönünde (birim: $\frac{1}{\sqrt5}\langle1,2\rangle$); en büyük hız $|\nabla f|=\sqrt5$.

\textbf{(c)} $\mathbf u\perp\langle1,2\rangle$: $\mathbf u=\pm\frac{1}{\sqrt5}\langle2,-1\rangle$.
\end{cevap}

\begin{soru}[\ekal Ek alıştırma 8 (üç değişken)]
$f(x,y,z)=xy+yz+zx$ için $P(1,-1,2)$ noktasında (a) $Q(2,1,4)$'e doğru yönlü türevi, (b) maksimum yönlü türevi bulunuz.
\end{soru}
\begin{cevap}
$\nabla f=\langle y+z,\ x+z,\ x+y\rangle$, $\nabla f(P)=\langle1,3,0\rangle$.

\textbf{(a)} $\overrightarrow{PQ}=\langle1,2,2\rangle$, uzunluğu $3$; $D_{\mathbf u}f=\frac{1\cdot1+3\cdot2+0\cdot2}{3}=\boxed{\frac73}$.

\textbf{(b)} $\max D_{\mathbf u}f=|\nabla f(P)|=\sqrt{1+9+0}=\boxed{\sqrt{10}}$, $\langle1,3,0\rangle$ yönünde.
\end{cevap}

\begin{soru}[\ekal Ek alıştırma 9 (teğet düzlem ve normal doğru)]
$x^2+y^2-z^2=4$ tek kanatlı hiperboloidinin $(2,1,1)$ noktasındaki teğet düzlemini ve normal doğrusunu bulunuz.
\end{soru}
\begin{cevap}
Kontrol: $4+1-1=4$ \checkmark. $F=x^2+y^2-z^2$, $\nabla F=\langle2x,2y,-2z\rangle$, $\nabla F(2,1,1)=\langle4,2,-2\rangle\parallel\langle2,1,-1\rangle$.

Teğet düzlem: $2(x-2)+(y-1)-(z-1)=0\iff\boxed{2x+y-z=4}$.

Normal doğru: $x=2+2t,\quad y=1+t,\quad z=1-t$ \ (simetrik: $\frac{x-2}{2}=\frac{y-1}{1}=\frac{z-1}{-1}$).
\end{cevap}
